\scanpage{47}

% 1
\[
 \alpha_{*}(\mathbb E^i)\simeq DP(C_{\bullet}^{i})\simeq\overset{i+1}{\bigwedge}\Phi_{*}
\]

% 2
\[
 (\text{d'où}\quad\alpha^{*}(\mathbb E^i)\simeq DP(C_{i}^{\bullet})\simeq\overset{i+1}{\bigwedge}\check{\Phi}_{*}\;)
\]

% 3–5, cancelled at upper edge
\deleted{L'\unsure{équ.|iso} déf. entre les $\alpha_{*}(\mathbb E^i)$ est l'\unsure{équiv. des} $T_{*}\wedge$ \unsure{que} le foncteur $ND:\mathbb K_{*}\longrightarrow\mathbb K_{\bullet}$ est \unsure{donné} \unsure{comme} de $(\overset{\unsure{\bullet\geq1|i+1}}{\bigwedge})\Phi_{*}$.}

% 5–7, crossed heading and formula
\deleted{Par suite \uline{Cor. 1} Le foncteur $DP:Ab_{k\,\bullet}\longrightarrow\unsure{Ab_{k\,*}}$}

% 8–11
Appliquant le \unsure{triple|théorie} de Dold-Puppe à $K$ \deleted{(qui est bien additive et stable par facteurs directs, du moins si $k$ est principal, p. ex $k=\mathbb Z$)} on trouve dans $K$ un complexe

% 12
\[
 \deleted{\unsure{C^{*}}}\qquad C^{*}=DP(C^{\bullet})
\]

% 13–14
\deleted{et \unsure{par|dans} Yoneda, pour} (qui est donc défini dans \deleted{\unsure{$K_k$}} $K(k)$ quel que soit $k$), d'où pour tte $B$ \inserted{$k$-}additive

% 15–18, triangle with intermediate restriction
\diagram{\begin{tikzcd}[ampersand replacement=\&,row sep=large]
 \underline{\operatorname{Hom}}_{k\,\mathrm{lin}}(\mathbb K,B)\arrow[d,"\mathrm{restr}","\wr"']\arrow[ddr,bend left=15,"{F\mapsto F(C^{*})}"]\&\\
 \underline{\operatorname{Hom}}_{k\text{-}\mathrm{lin}}(\unsure{\mathscr A|\mathcal A},B)\arrow[d,"{F\mapsto F(C^{\bullet})}","\wr"']\&\\
 B^{\bullet}\arrow[r,leftrightarrow,"{DP_{\unsure{\bullet|\,}}}"]\& B^{*}
\end{tikzcd}}

% 19–23
donc un foncteur $F\mapsto F(C^{*})$ qui est pl. fid. Si $B$ est karoubienne, $DP_B$ est une équivalence, donc aussi $F\mapsto F(C^{*})$. Donc $C^{*}$ apparaît comme objet \unsure{cosimplicial|simplicial} \uline{universel} dans cat. $k$-Additive karoubienne

% 24–25
\uline{Prop 5} Si $B$ est $k$-additive karoubienne, alors $F\mapsto F(C^{*})$ est une équivalence

% 26
\[
 \underline{\operatorname{Hom}}_{k\text{-}\mathrm{lin}}(\mathbb K,B)\longrightarrow B^{*}
\]

% 27–28
On a un isom canonique \inserted{bifonct.l} en $F\in\operatorname{Ob}\underline{\operatorname{Hom}}_{k\text{-}\mathrm{lin}}(\mathbb K,B)$ et $x\in\operatorname{Ob}\mathbb K$.

% 28, lower edge below Hom
\margin{\unsure{$\wr$}}

% archive, bottom left
% archive sheet number: 46
